\documentclass[UTF8, a4paper, 12pt]{ctexart}
\usepackage{amsmath, amssymb, booktabs, geometry}
\geometry{left=2.5cm, right=2.5cm, top=3cm, bottom=3cm}
\usepackage{color}



\title{\textbf{Timoshenko 樑 (Timoshenko Beam) 有限元素分析與驗證報告}}
\author{Finite Element Analysis Team}
\date{\today}

\begin{document}
\maketitle

\section{問題描述與參數設定 (Problem Description \& Parameters)}

本報告針對 Timoshenko 樑（含剪力形變與旋轉慣量效應）進行有限元素分析，以驗證靜態彎曲、屈曲與自由震動的正確性。根據 MATLAB Codes for Finite Element Analysis (Ferreira 2009) 之設定，本模型採用以下基本物理與幾何參數：

\begin{itemize}
  \item \textbf{楊氏模數 (Young's modulus, $E$)}: $210 \text{ GPa}$
  \item \textbf{波松比 (Poisson's ratio, $\nu$)}: $0.3$
  \item \textbf{剪力模數 (Shear modulus, $G$)}: $G = E / 2(1+\nu) \approx 80.77 \text{ GPa}$
  \item \textbf{樑長度 (Length, $L$)}: $10 \text{ m}$
  \item \textbf{截面面積 (Area, $A$)}: $0.02 \text{ m}^2$
  \item \textbf{慣性矩 (Moment of inertia, $I$)}: $10 \times 10^{-5} \text{ m}^4$
  \item \textbf{剪力修正係數 (Shear correction factor, $\kappa$)}: $5/6$
  \item \textbf{密度 (Density, $\rho$)}: $1.0\text{ (歸一化參數供頻率比對)}$
  \item \textbf{均布載重 (Uniform load, $q$)}: $1000 \text{ N/m}$ (用於靜態分析)
  \item \textbf{網格劃分 (Elements)}: 三個模型均採用 $20\sim100$ 個元素
\end{itemize}

\section{理論基礎與數值方法對比 (Methodology Comparison)}

本研究將解析理論解與有限元素法 (FEM) 進行對比。Timoshenko 樑理論與傳統 Euler-Bernoulli 理論的最大差異在於考慮了截面的剪力形變 ($\gamma$) 與旋轉慣量效應。

\subsection{Timoshenko 樑理論完整推導}

以下為 Timoshenko 樑理論的完整數學推導過程，採用無鎖定（Locking-Free）有限元素公式。

%\subsubsection{1.1 鐵摩辛柯梁理論基本概念}

鐵摩辛柯樑理論是對歐拉–伯努力樑（Euler–Bernoulli beam）理論的延伸，能同時考慮剪力變形與截面轉動所造成的彎曲效應。它對厚樑、或高頻振動等情況特別重要。

其運動學假設為：截面在變形後仍保持平面，但不一定仍與樑的縱向軸線保持垂直。控制方程涉及兩個相互獨立的未知量：橫向位移 $w(x)$ 與截面轉角 $\phi(x)$。

%\subsubsection{1.2 剪力鎖定（Shear Locking）現象}

剪力鎖定是有限元素分析中的一種數值人為效應。當對具有顯著剪力變形的問題使用標準低階元素時，常會出現此現象。

\textbf{成因：} 當樑非常細長（薄）時，元素難以正確表示「純彎曲」狀態；剪應變被人為地限制為非零，導致結構回應出現「鎖住」且過度剛硬（overly stiff）。

\textbf{影響：} 隨著網格細化，計算得到的位移會趨近於零，而非正確解；對薄樑尤其嚴重。

%\subsubsection{1.3 無鎖定（Locking-Free）公式}

為了克服剪力鎖定，已發展出多種進階公式。本專案採用：
\begin{itemize}
  \item \textbf{降階積分（Reduced Integration）：} 對元素剛度矩陣中的剪力項採用較低階的積分規則。
  \item \textbf{混合式公式（Mixed Formulations）：} 將位移場與應力／應變場視為彼此獨立的未知量。
\end{itemize}

%\subsubsection{1.4 屈曲分析的位能泛函}

屈曲分析可由「位能駐值原理」（principle of stationary potential energy）出發。樑在軸向壓縮力 $P$ 作用下，其總位能 $\Pi$ 由應變能 $U$ 與外力位能 $V$ 所構成。
\begin{equation}
  \Pi = U - V
\end{equation}

\textbf{彎曲應變能（$U_{b}$）：} 由彎曲造成的儲能。
\begin{equation}
  U_{b} = \frac{1}{2}\int_{0}^{L}EI\phi_{,x}^{2}dx
\end{equation}
其中 $E$ 為楊氏模數，$I$ 為截面二次矩，$\phi$ 為截面轉角。

\textbf{剪力應變能（$U_s$）：} 由剪力變形造成的儲能。
\begin{equation}
  U_{s} = \frac{1}{2}\int_{0}^{L}kAG(w_{,x}-\phi)^{2}dx
\end{equation}
其中，$k$ 為剪力修正係數，$A$ 為截面面積，$G$ 為剪力模數，$w$ 為橫向位移。

\textbf{外力位能：} 軸向壓縮力 $P$ 的位能可視為該力因樑的橫向變形所做的功：
\begin{equation}
  V = \frac{1}{2}\int_{0}^{L}Pw_{,x}^{2}dx
\end{equation}

\textbf{總位能泛函：}
\begin{equation}
  \Pi = \frac{1}{2}\int_{0}^{L}(EI\phi_{,x}^{2}+kAG(w_{,x}-\phi)^{2}-Pw_{,x}^{2})dx
\end{equation}

%\subsubsection{1.5 變分原理與控制方程}

平衡狀態可由位能駐值原理求得：對任何允許的虛位移，總位能的一次變分必須為零。
\begin{equation}
  \delta\Pi = \delta\frac{1}{2}\int_{0}^{L}(EI\phi_{,x}^{2}+kAG(w_{,x}-\phi)^{2}-Pw_{,x}^{2})dx = 0
\end{equation}

對獨立未知量 $w$ 與 $\phi$ 取變分可得：

\textbf{對 $w$ 的弱式：}
\begin{equation}
  \int_{0}^{L}(\delta w_{,x}kAGw_{,x} - \delta w_{,x}kAG\phi - \delta w_{,x}Pw_{,x})dx = 0
\end{equation}

\textbf{對 $\phi$ 的弱式：}
\begin{equation}
  \int_{0}^{L}(\delta\phi_{,x}EI\phi_{,x} - \delta\phi kAGw_{,x} + \delta\phi kAG\phi)dx = 0
\end{equation}

%\subsubsection{1.6 有限元素近似}

將樑域離散成有限元素，並以形狀函數近似場變量 $w$ 與 $\phi$。
\begin{equation}
  w \approx w^{h} = \sum_{I=1}^{n_{w}}N_{I}(x)w_{I}
\end{equation}
\begin{equation}
  \phi \approx \phi^{h} = \sum_{J=1}^{n_{\phi}}N_{J}(x)\phi_{J}
\end{equation}

其導數近似為：
\begin{equation}
  w_{,x} \approx w_{,x}^{h} = \sum_{I=1}^{n_{w}}N_{I,x}(x)w_{I}
\end{equation}
\begin{equation}
  \phi_{,x} \approx \phi_{,x}^{h} = \sum_{J=1}^{n_{\phi}}N_{J,x}(x)\phi_{J}
\end{equation}

%\subsubsection{1.7 離散後的特徵值問題}

將上述近似代入變分式，即可得到離散後的特徵值問題：
\begin{equation}
  \begin{bmatrix}
    K_{ww} & K_{w\phi} \\
    K_{\phi w} & K_{\phi\phi}
  \end{bmatrix}
  \begin{bmatrix}
    d_{w} \\
    d_{\phi}
  \end{bmatrix}
  - P
  \begin{bmatrix}
    K_{,ww}^{G} & 0 \\
    0 & 0
  \end{bmatrix}
  \begin{bmatrix}
    d_{w} \\
    d_{\phi}
  \end{bmatrix} = \begin{bmatrix} 0 \\ 0 \end{bmatrix}
\end{equation}

其中各元素剛度子矩陣定義為：
\begin{equation}
  K_{ww,IJ} = \int_{0}^{L}(kAGN_{I,x}N_{J,x} - PN_{I,x}N_{J,x})dx
\end{equation}
\begin{equation}
  K_{\phi\phi,IJ} = \int_{0}^{L}(EIN_{I,x}N_{J,x} + kAGN_{I}N_{J})dx
\end{equation}
\begin{equation}
  K_{w\phi,IJ} = K_{\phi w,JI} = -\int_{0}^{L}kAGN_{I,x}N_{J}dx
\end{equation}
\begin{equation}
  K_{,ww,IJ}^{G} = \int_{0}^{L}N_{I,x}N_{J,x}dx
\end{equation}

全域剛度矩陣由各元素貢獻組裝而成。臨界屈曲荷載 $P_{cr}$ 是使得總位能的二次變分不再為正定的最小 $P$，因此會導出一個特徵值問題。

\subsection{靜態彎曲位移公式推導}

本節詳細說明 Timoshenko 樑靜態彎曲分析中，總位移公式 $w = w_b + w_s$ 的推導過程及其線性性質。

\paragraph{2.2.1 Timoshenko 樑理論的基本假設}

Timoshenko 樑理論放寬了 Euler-Bernoulli 理論的「平面保持平面且垂直於中性軸」假設：

\begin{itemize}
  \item 截面在變形後仍保持平面，但不一定與中性軸垂直
  \item 位移場定義為：
  \begin{itemize}
    \item 軸向位移：$u(x, y) = y \cdot \theta(x)$
    \item 橫向位移：$w(x, y) = w_0(x)$
  \end{itemize}
\end{itemize}

其中 $\theta(x)$ 是截面旋轉角度，$w_0(x)$ 是中性軸的橫向位移。

\paragraph{2.2.2 應變分析}

由位移場可求得應變：

\begin{itemize}
  \item \textbf{彎曲應變（正應變）：}
  $$\varepsilon_x = \frac{\partial u}{\partial x} = y \frac{d\theta}{dx}$$
  
  \item \textbf{剪力應變：}
  $$\gamma_{xy} = \frac{\partial u}{\partial y} + \frac{\partial w}{\partial x} = \theta + \frac{dw_0}{dx}$$
\end{itemize}

\paragraph{2.2.3 應變能}

總應變能包含彎曲與剪力兩部分。

$$U = \underbrace{\frac{1}{2}\int_0^L EI\left(\frac{d\theta}{dx}\right)^2 dx}_{U_b} + \underbrace{\frac{1}{2}\int_0^L \kappa AG\left(\frac{dw_0}{dx} - \theta\right)^2 dx}_{U_s}$$

\paragraph{2.2.4 位移公式推導}

由變分原理可得控制方程。對於均布載重 $q$ 作用下的樑，總位移為：

$$w = w_b + w_s$$

其中：
\begin{itemize}
  \item \textbf{彎曲貢獻 ($w_b$)：} 來源於彎曲應變能 $U_b$，與 Euler-Bernoulli 樑理論相同。
  \begin{equation}
    w_b = \frac{qL^4}{8EI} \quad \text{(懸臂樑)} \quad \text{或} \quad w_b = \frac{5qL^4}{384EI} \quad \text{(簡支樑)}
  \end{equation}
  
  \item \textbf{剪切貢獻 ($w_s$)：} 來源於剪力應變能 $U_s$，這是 Timoshenko 理論特有的項。
  \begin{equation}
    w_s = \frac{qL^2}{2\kappa GA} \quad \text{(懸臂樑)} \quad \text{或} \quad w_s = \frac{qL^2}{8\kappa GA} \quad \text{(簡支樑)}
  \end{equation}
\end{itemize}

\textbf{重要結論：} $w = w_b + w_s$ 是\textbf{線性方程}。

證明：在小變形、線彈性假設下（虎克定律成立），位移與載重成正比：

$$w_b = \frac{qL^4}{8EI} \propto q \quad \text{（一次方）}$$
$$w_s = \frac{qL^2}{2\kappa GA} \propto q \quad \text{（一次方）}$$

因此，
$$w = w_b + w_s \propto q$$

這稱為\textbf{疊加原理}——在線性系統中，總位移等於各單獨載重作用下位移的代數和。

\subsection{自由震動分析的解析解推導}

本節詳細說明 Timoshenko 樑自由震動分析中，自然頻率的推導過程。與靜態分析不同，震動分析需要考慮慣性效應，包括平移慣性和旋轉慣性。

\subparagraph{2.3.1 運動方程的推導}

對於自由震動（無外力作用），樑的運動由拉格朗日方程描述。總動能 $T$ 包含兩部分：

\begin{itemize}
  \item \textbf{平移動能 ($T_w$)：} 來自橫向位移 $w(x,t)$
  \begin{equation}
    T_w = \frac{1}{2}\int_0^L \rho A \dot{w}^2 \, dx
  \end{equation}
  
  \item \textbf{旋轉動能 ($T_\phi$)：} 來自截面轉角 $\phi(x,t)$
  \begin{equation}
    T_\phi = \frac{1}{2}\int_0^L \rho I \dot{\phi}^2 \, dx
  \end{equation}
\end{itemize}

其中，$\dot{w} = \partial w/\partial t$，$\dot{\phi} = \partial \phi/\partial t$ 為時間導數，$\rho$ 為密度，$I$ 為慣性矩。

\subparagraph{2.3.2 震動控制方程}

總位能 $U$ 與靜態分析相同：
\begin{equation}
  U = \frac{1}{2}\int_0^L \left[ EI\left(\frac{\partial \phi}{\partial x}\right)^2 + \kappa GA\left(\frac{\partial w}{\partial x} - \phi\right)^2 \right] dx
\end{equation}

根據拉格朗日方程，對於廣義坐標 $w$ 和 $\phi$：
\begin{equation}
  \frac{d}{dt}\left(\frac{\partial T}{\partial \dot{w}}\right) - \frac{\partial T}{\partial w} + \frac{\partial U}{\partial w} = 0
\end{equation}
\begin{equation}
  \frac{d}{dt}\left(\frac{\partial T}{\partial \dot{\phi}}\right) - \frac{\partial T}{\partial \phi} + \frac{\partial U}{\partial \phi} = 0
\end{equation}

代入動能和位能表達式，可得 Timoshenko 樑的自由震動控制方程：
\begin{equation}
  \kappa GA \left( \frac{\partial^2 w}{\partial x^2} - \frac{\partial \phi}{\partial x} \right) = \rho A \frac{\partial^2 w}{\partial t^2}
\end{equation}
\begin{equation}
  EI \frac{\partial^2 \phi}{\partial x^2} + \kappa GA \left( \frac{\partial w}{\partial x} - \phi \right) = \rho I \frac{\partial^2 \phi}{\partial t^2}
\end{equation}

\subparagraph{2.3.3 頻率方程的推導}

對於簡諧震動，假設解的形式為：
\begin{equation}
  w(x,t) = W(x) e^{i\omega t}, \quad \phi(x,t) = \Phi(x) e^{i\omega t}
\end{equation}
其中 $\omega$ 為圓頻率。

代入控制方程，得到特徵值問題：
\begin{equation}
  \kappa GA \left( W'' - \Phi' \right) + \rho A \omega^2 W = 0
\end{equation}
\begin{equation}
  EI \Phi'' + \kappa GA (W' - \Phi) + \rho I \omega^2 \Phi = 0
\end{equation}

經過數學整理，可得 Timoshenko 樑的頻率方程。對於懸臂樑（固定端 $x=0$，自由端 $x=L$），特徵值 $\beta$ 滿足：
\begin{equation}
  \beta^4 - \alpha \beta^2 - \gamma = 0
\end{equation}
其中：
\begin{equation}
  \alpha = \frac{\rho A L^4}{EI}, \quad \gamma = \frac{\rho I L^2}{EI}
\end{equation}

\subparagraph{2.3.4 自然頻率公式}

懸臂樑的前四階自然頻率為：
\begin{equation}
  \omega_n = \beta_n^2 \sqrt{\frac{EI}{\rho A L^4}} \cdot \sqrt{\frac{1}{1 + \gamma \beta_n^2}}
\end{equation}

簡支樑的自然頻率為：
\begin{equation}
  \omega_n = \beta_n^2 \sqrt{\frac{EI}{\rho A L^4}} \cdot \sqrt{\frac{1}{1 + (\gamma + \alpha)\beta_n^2}}
\end{equation}

其中 $\beta_n = n\pi$（簡支樑）或由超越方程決定（懸臂樑）。

\textbf{重要結論：} Timoshenko 理論考慮了旋轉慣量效應，導致頻率低於 Euler-Bernoulli 理論的預測值。這個效應在高頻模態和厚樑中特別顯著。

\subsection{自由震動分析的有限元素法推導}

本節說明如何將自由震動問題離散化為有限元素形式。

\subparagraph{2.4.1 質量矩陣的推導}

使用與靜態分析相同的形狀函數近似場變量：
\begin{equation}
  w \approx w^h = \sum_{I=1}^{n} N_I(x) w_I(t), \quad \phi \approx \phi^h = \sum_{I=1}^{n} N_I(x) \phi_I(t)
\end{equation}

動能離散化為：
\begin{equation}
  T = \frac{1}{2} \dot{d}^T M \dot{d}
\end{equation}

其中 $\dot{d} = [\dot{w}_1, \dot{\phi}_1, \dots, \dot{w}_n, \dot{\phi}_n]^T$ 為廣義速度向量，$M$ 為總質量矩陣。

\subparagraph{2.4.2 元素質量矩陣}

對於兩節點 Timoshenko 樑元素，質量矩陣為：
\begin{equation}
  M_e = 
  \begin{bmatrix}
    M_{ww} & M_{w\phi} \\
    M_{\phi w} & M_{\phi\phi}
  \end{bmatrix}
\end{equation}

其中各子矩陣定義為：
\begin{equation}
  M_{ww,IJ} = \int_0^{l_e} \rho A N_I N_J \, dx = \frac{\rho A l_e}{6}
  \begin{bmatrix}
    2 & 0 \\
    0 & 2 \\
    1 & 0 \\
    0 & 1
  \end{bmatrix}
\end{equation}
\begin{equation}
  M_{\phi\phi,IJ} = \int_0^{l_e} \rho I N_I N_J \, dx = \frac{\rho I l_e}{6}
  \begin{bmatrix}
    2 & 0 \\
    0 & 2 \\
    1 & 0 \\
    0 & 1
  \end{bmatrix}
\end{equation}
\begin{equation}
  M_{w\phi,IJ} = M_{\phi w,JI} = 0
\end{equation}

\textbf{重要特性：} Timoshenko 樑的質量矩陣包含兩個獨立的質量項：
\begin{itemize}
  \item 平移質量矩陣：$\rho A$（與桿件質量相關）
  \item 旋轉質量矩陣：$\rho I$（與截面轉動慣量相關）
\end{itemize}

旋轉慣量項 $\rho I$ 是 Timoshenko 理論與 Euler-Bernoulli 理論的本質區別。

\subparagraph{2.4.3 離散後的特徵值問題}

將位能離散化後，總質量矩陣組裝為全域矩陣 $M$。自由震動問題的平衡方程為：
\begin{equation}
  M \ddot{d} + K d = 0
\end{equation}

代入諧波假設 $d(t) = d_0 e^{i\omega t}$，得到標準特徵值問題：
\begin{equation}
  (K - \omega^2 M) \{d\} = \{0\}
\end{equation}

求解此特徵值問題即可得到系統的自然頻率 $\omega$ 和振型 $\{d\}$。

\subparagraph{2.4.4 FEM 與解析解的對比}

有限元素法通過求解廣義特徵值問題獲得頻率，而解析法則直接求解超越方程。兩種方法的對比：

\begin{center}
\begin{tabular}{p{3cm}p{5cm}p{5cm}}
\toprule
\textbf{特性} & \textbf{解析法} & \textbf{有限元素法} \\
\midrule
\textbf{數學本質} & 超越方程求解 & 矩陣特徵值求解 \\
\midrule
\textbf{網格依賴} & 無 & 網格細化收斂 \\
\midrule
\textbf{計算效率} & 單一模態快速 & 一次計算多模態 \\
\midrule
\textbf{幾何複雜度} & 僅限簡單邊界 & 可處理複雜幾何 \\
\bottomrule
\end{tabular}
\end{center}

本報告採用 $20\sim100$ 個元素進行收斂性分析，驗證 FEM 結果的正確性。

\subsection{ 核心公式運用 (Timoshenko Formulas)}

下表彙整了本專案中使用的兩種計算途徑：

\begin{center}
\begin{tabular}{p{3cm}p{6cm}p{6cm}}
\toprule
\textbf{特性} & \textbf{解析計算 (Analytical)} & \textbf{本專案測試方法 (FEM)} \\
\midrule
\textbf{數學本質} & 封閉解公式 (Closed-form) & 矩陣位移法與特徵值求解 \\
\midrule
\textbf{剪切處理} & 引入剪切修正因子 (Shear factor) & 採用 \textbf{單點高斯積分} (1-point integration) 消除剪力鎖定 \\
\midrule
  \textbf{空間離散} & 連續函數 & 將樑劃分為 20 ~ 100 個元素  \\
\midrule
\textbf{求解器} & 直接帶入公式 & 調用 Julia \texttt{LinearAlgebra} 與 \texttt{eigvals} 函數 \\
\bottomrule
\end{tabular}
\end{center}

\subsection{ 程式碼實現說明}

根據本專案開發的 Julia 檔案，以下為 Timoshenko 理論在程式碼中的具體實作方式：

\begin{itemize}
  \item \textbf{剛度矩陣構造：}
  在 \texttt{TimoshenkoBeam.jl} 中，剛度矩陣被拆分為彎曲部分 $K_b$ 與剪力部分 $K_s$。剪力部分使用了 \texttt{(κGA * le) * (B\_s * B\_s')}，這正是 Timoshenko 理論中剪力剛度的數值呈現。
  
  \item \textbf{消除剪力鎖定 (Shear Locking)：}
  程式碼中特別採用「單點積分」(1-point integration)。這是 Timoshenko 樑有限元素分析中的關鍵技術，若使用標準積分會導致樑在變薄時剛度過大，無法得到準確的解析對比值。

\subparagraph{單點積分與完整積分的比較}

本專案同時實現了兩種積分方式，用於對比分析：

\textbf{1. 單點積分（Reduced Integration / 1-point Gauss）：}
\begin{itemize}
  \item 對剪力剛度矩陣使用 1 個高斯點進行數值積分
  \item 高斯點位置：$\xi = 0$（元素中心）
  \item 權重：$w = 2$
  \item \textbf{優點：} 有效消除剪力鎖定問題，計算效率高
  \item \textbf{缺點：} 可能導致零能量模式（hourglassing），但可通過適當的網格設計避免
\end{itemize}

\textbf{2. 完整積分（Full Integration / 2-point Gauss）：}
\begin{itemize}
  \item 對剪力剛度矩陣使用 2 個高斯點進行數值積分
  \item 高斯點位置：$\xi = \pm 1/\sqrt{3}$
  \item 權重：$w = 1$（每個點）
  \item \textbf{優點：} 精確積分，無零能量模式
  \item \textbf{缺點：} 會產生剪力鎖定問題，導致結構過度剛硬
\end{itemize}

\textbf{為什麼 2 個高斯點是「完整」積分？}

對於兩節點 Timoshenko 樑元素：
\begin{itemize}
  \item 位移場使用線性形狀函數（1階多項式）
  \item 剪切應變矩陣 $B_s$ 包含形狀函數的導數，為常數（0階）
  \item 被積分函數為 2 階多項式（兩個常數相乘）
\end{itemize}

要精確積分 $n$ 階多項式，需要至少 $n/2 + 1$ 個高斯點。因此對於 2 階多項式，2 個高斯點就是完整積分。

\textbf{數學表達：}
\begin{equation}
  \text{單點積分：} \quad \int_{-1}^{1} f(\xi) d\xi \approx f(0) \cdot 2
\end{equation}
\begin{equation}
  \text{完整積分：} \quad \int_{-1}^{1} f(\xi) d\xi \approx \sum_{i=1}^{2} f(\xi_i) \cdot w_i, \quad \xi_i = \pm\frac{1}{\sqrt{3}}, w_i = 1
\end{equation}

本專案採用單點積分版本的 \texttt{TimoshenkoBeam.jl} 作為主要實現，並提供完整積分版本 \texttt{TimoshenkoBeamFullIntegration.jl} 用於教學對比。

  \item \textbf{幾何剛度與質量矩陣：}
  \begin{itemize}
    \item \textbf{屈曲}：\texttt{TimoshenkoBeamBuckling.jl} 構造了幾何剛度矩陣 \texttt{KG\_local}，用來捕捉軸向力對橫向剛度的影響。
    \item \textbf{震動}：\texttt{TimoshenkoBeamVibrations.jl} 在質量矩陣 \texttt{M\_local} 中同時包含了平移質量 (\texttt{rho * A}) 與旋轉質量 (\texttt{rho * I})，這完全符合 Timoshenko 對高頻模態的描述。
  \end{itemize}
\end{itemize}

\section{有限元素法結果與解析解對比 (Results Comparison)}

以下展示有限元素法（FEM）網格收斂性分析，元素數量由 20 至 100，比較 FEM 結果與解析解的誤差。

\subsection{靜態彎曲分析 (Static Bending Analysis)}

\textbf{解析解公式：}
彎曲位移與剪力位移之和為 Timoshenko 樑的總位移：
\begin{equation}
  w = w_b + w_s
\end{equation}

其中彎曲貢獻為：
\begin{equation}
  w_b = \frac{qL^4}{8EI} \quad \text{(懸臂樑)} \quad \text{或} \quad w_b = \frac{5qL^4}{384EI} \quad \text{(簡支樑)}
\end{equation}

剪力貢獻為：
\begin{equation}
  w_s = \frac{qL^2}{2\kappa GA} \quad \text{(懸臂樑)} \quad \text{或} \quad w_s = \frac{qL^2}{8\kappa GA} \quad \text{(簡支樑)}
\end{equation}

\textbf{FEM 代碼實現：}
\begin{itemize}
  \item \textbf{彎曲剛度矩陣：}
  \begin{equation}
    K_b = \frac{EI}{l_e} \begin{bmatrix} 0 & 0 & 0 & 0 \\ 0 & 1 & 0 & -1 \\ 0 & 0 & 0 & 0 \\ 0 & -1 & 0 & 1 \end{bmatrix}
  \end{equation}
  \item \textbf{剪力剛度矩陣（1-point 積分）：}
  \begin{equation}
    K_s = \kappa GA \cdot l_e \cdot B_s B_s^T, \quad B_s = \begin{bmatrix} -\frac{1}{l_e} & -\frac{1}{2} & \frac{1}{l_e} & -\frac{1}{2} \end{bmatrix}
  \end{equation}
  \item \textbf{總剛度矩陣：} $K = K_b + K_s$
\end{itemize}

\paragraph{(a) 簡支樑 (Simply-supported)}

\textbf{收斂性分析結果：}

\begin{center}
\begin{tabular}{cccc}
\toprule
\textbf{元素數} & \textbf{FEM (m)} & \textbf{解析解 (m)} & \textbf{誤差 (\%)} \\
\midrule
20 & 0.00618488 & 0.00620968 & 0.3994 \\
30 & 0.00619866 & 0.00620968 & 0.1775 \\
40 & 0.00620348 & 0.00620968 & 0.0999 \\
50 & 0.00620571 & 0.00620968 & 0.0640 \\
100 & 0.00620869 & 0.00620968 & 0.0160 \\
\bottomrule
\end{tabular}
\end{center}

\textit{注：隨著元素數量增加，位移誤差持續下降，展現良好的收斂性。}

\paragraph{(b) 懸臂樑 (Cantilever)}

\textbf{收斂性分析結果：}

\begin{center}
\begin{tabular}{cccc}
\toprule
\textbf{元素數} & \textbf{FEM (m)} & \textbf{解析解 (m)} & \textbf{誤差 (\%)} \\
\midrule
20 & 0.0595609524 & 0.0595609524 & $0$ \\
30 & 0.0595609524 & 0.0595609524 & $0$ \\
40 & 0.0595609524 & 0.0595609524 & $0$ \\
50 & 0.0595609524 & 0.0595609524 & $0$ \\
100 & 0.0595609524 & 0.0595609524 & $0$ \\
\bottomrule
\end{tabular}
\end{center}

\textit{注：解析解與 FEM 結果完全吻合，驗證了程式實現的正確性。}

\subparagraph{(c) 完整積分（Full Integration）對比}

本節展示使用完整積分（2-point Gauss）時的 FEM 結果，與解析解進行對比，以說明剪力鎖定問題的影響。完整積分使用 2 個高斯點（$\xi = \pm 1/\sqrt{3}$）進行數值積分。

\textbf{簡支樑收斂性分析結果：}

\begin{center}
\begin{tabular}{cccc}
\toprule
\textbf{元素數} & \textbf{FEM (m)} & \textbf{解析解 (m)} & \textbf{誤差 (\%)} \\
\midrule
20 & 0.00168694 & 0.00620968 & 72.8338 \\
30 & 0.00283461 & 0.00620968 & 54.3518 \\
40 & 0.00371878 & 0.00620968 & 40.1132 \\
50 & 0.00434585 & 0.00620968 & 30.0149 \\
100 & 0.00560565 & 0.00620968 & 9.7273 \\
\bottomrule
\end{tabular}
\end{center}

\textbf{懸臂樑收斂性分析結果：}

\begin{center}
\begin{tabular}{cccc}
\toprule
\textbf{元素數} & \textbf{FEM (m)} & \textbf{解析解 (m)} & \textbf{誤差 (\%)} \\
\midrule
20 & 0.01623344 & 0.05956095 & 72.7448 \\
30 & 0.02723463 & 0.05956095 & 54.2744 \\
40 & 0.03570998 & 0.05956095 & 40.0446 \\
50 & 0.04172088 & 0.05956095 & 29.9526 \\
100 & 0.05379684 & 0.05956095 & 9.6777 \\
\bottomrule
\end{tabular}
\end{center}

\textit{注：完整積分（2-point Gauss）導致嚴重的剪力鎖定問題。即使增加到 100 個元素，誤差仍有約 10\%，無法像單點積分那樣收斂到精確解。這是因為完整積分高估了結構的剪切剛度，導致位移被嚴重低估。}

\subparagraph{單點積分與完整積分對比總結}

\begin{center}
\begin{tabular}{p{2.5cm}p{4cm}p{4cm}}
\toprule
\textbf{項目} & \textbf{單點積分 (1-point)} & \textbf{完整積分 (2-point)} \\
\midrule
高斯點數量 & 1 個 ($\xi = 0$) & 2 個 ($\xi = \pm 1/\sqrt{3}$) \\
\midrule
簡支樑誤差 (100 元素) & 0.0160\% & 9.7273\% \\
\midrule
懸臂樑誤差 (100 元素) & 0\% & 9.6777\% \\
\midrule
收斂特性 & 良好收斂 & 收斂緩慢，無法完全收斂 \\
\midrule
剪力鎖定 & 無 & 嚴重 \\
\bottomrule
\end{tabular}
\end{center}

\textbf{結論：} 單點積分有效消除剪力鎖定問題，是 Timoshenko 樑有限元素分析的首選數值積分方法。完整積分雖然看起來更「精確」，但實際上會導致嚴重的數值問題。

\subsection{屈曲分析 (Buckling Analysis)}

\textbf{解析解公式：}
屈曲分析的臨界載重由 Euler 公式經 Timoshenko 修正得到。
\begin{equation}
  P_{cr} = \frac{P_{cr}^{Euler}}{1 + \frac{P_{cr}^{Euler}}{\kappa AG}}
\end{equation}

其中 Euler 臨界載重為：
\begin{equation}
  P_{cr}^{Euler} = \frac{\pi^2 EI}{L^2} \quad \text{(雙端鉸接)} \quad \text{或} \quad P_{cr}^{Euler} = \frac{\pi^2 EI}{4L^2} \quad \text{(懸臂樑)}
\end{equation}

\textbf{FEM 代碼實現：}
\begin{itemize}
  \item \textbf{幾何剛度矩陣：}
  \begin{equation}
    K_G = \frac{1}{l_e} \begin{bmatrix} 1 & 0 & -1 & 0 \\ 0 & 0 & 0 & 0 \\ -1 & 0 & 1 & 0 \\ 0 & 0 & 0 & 0 \end{bmatrix}
  \end{equation}
  \item \textbf{特徵值問題：}
  \begin{equation}
    [K - P K_G] \{u\} = \{0\}
  \end{equation}
  其中 $P$ 為特徵值（臨界載重）。
\end{itemize}

\subparagraph{(a) 雙端鉸接 (Pinned-Pinned)}

\textbf{收斂性分析結果：}

\begin{center}
\begin{tabular}{cccc}
\toprule
\textbf{元素數} & \textbf{FEM (N)} & \textbf{解析解 (N)} & \textbf{誤差 (\%)} \\
\midrule
20 & $2.078\times10^6$ & $2.069\times10^6$ & 0.4120 \\
30 & $2.073\times10^6$ & $2.069\times10^6$ & 0.1828 \\
40 & $2.072\times10^6$ & $2.069\times10^6$ & 0.1027 \\
50 & $2.071\times10^6$ & $2.069\times10^6$ & 0.0657 \\
100 & $2.070\times10^6$ & $2.069\times10^6$ & 0.0164 \\
\bottomrule
\end{tabular}
\end{center}

\subparagraph{(b) 懸臂樑 (Cantilever)}

\textbf{收斂性分析結果：}

\begin{center}
\begin{tabular}{cccc}
\toprule
\textbf{元素數} & \textbf{FEM (N)} & \textbf{解析解 (N)} & \textbf{誤差 (\%)} \\
\midrule
20 & $5.188\times10^5$ & $5.180\times10^5$ & 0.1029 \\
30 & $5.192\times10^5$ & $5.180\times10^5$ & 0.0457 \\
40 & $5.188\times10^5$ & $5.180\times10^5$ & 0.0257 \\
50 & $5.180\times10^5$ & $5.180\times10^5$ & 0.0164 \\
100 & $5.180\times10^5$ & $5.180\times10^5$ & 0.0041 \\
\bottomrule
\end{tabular}
\end{center}

\textit{注：屈曲臨界載重隨著網格細化趨近解析解，展現穩定收斂特性。}

\subsection{自由震動分析 (Free Vibrations)}

\textbf{解析解公式：}
自然頻率由 Euler-Bernoulli 理論給出：
\begin{equation}
  \omega_n = \beta_n^2 \sqrt{\frac{EI}{\rho A}} \cdot \frac{1}{L^2}
\end{equation}

其中 $\beta_n$ 為懸臂樑的特徵值：
\begin{equation}
  \beta = [1.8751,\ 4.6941,\ 7.8548,\ 10.9955,\dots]
\end{equation}

\textbf{FEM 代碼實現：}
\begin{itemize}
  \item \textbf{平移質量矩陣：}
  \begin{equation}
    M_w = \frac{\rho A \cdot l_e}{6} \begin{bmatrix} 2 & 0 & 1 & 0 \\ 0 & 0 & 0 & 0 \\ 1 & 0 & 2 & 0 \\ 0 & 0 & 0 & 0 \end{bmatrix}
  \end{equation}
  \item \textbf{旋轉質量矩陣：}
  \begin{equation}
    M_\theta = \frac{\rho I \cdot l_e}{6} \begin{bmatrix} 0 & 0 & 0 & 0 \\ 0 & 2 & 0 & 1 \\ 0 & 0 & 0 & 0 \\ 0 & 1 & 0 & 2 \end{bmatrix}
  \end{equation}
  \item \textbf{總質量矩陣：} $M = M_w + M_\theta$
  \item \textbf{特徵值問題：} $[K - \omega^2 M] \{u\} = \{0\}$
\end{itemize}

\subparagraph{Mode 1 收斂性分析：}

\begin{center}
\begin{tabular}{cccc}
\toprule
\textbf{元素數} & \textbf{FEM (rad/s)} & \textbf{解析解 (rad/s)} & \textbf{誤差 (\%)} \\
\midrule
20 & 1139.12 & 1139.31 & 0.0170 \\
30 & 1138.93 & 1139.31 & 0.0339 \\
40 & 1138.86 & 1139.31 & 0.0398 \\
50 & 1138.83 & 1139.31 & 0.0425 \\
100 & 1138.79 & 1139.31 & 0.0462 \\
\bottomrule
\end{tabular}
\end{center}

\subparagraph{Mode 2 收斂性分析：}

\begin{center}
\begin{tabular}{cccc}
\toprule
\textbf{元素數} & \textbf{FEM (rad/s)} & \textbf{解析解 (rad/s)} & \textbf{誤差 (\%)} \\
\midrule
20 & 7159.96 & 7140.02 & 0.2792 \\
30 & 7135.66 & 7140.02 & 0.0611 \\
40 & 7127.17 & 7140.02 & 0.1799 \\
50 & 7123.25 & 7140.02 & 0.2349 \\
100 & 7118.02 & 7140.02 & 0.3081 \\
\bottomrule
\end{tabular}
\end{center}

\subparagraph{Mode 3 收斂性分析：}

\begin{center}
\begin{tabular}{cccc}
\toprule
\textbf{元素數} & \textbf{FEM (rad/s)} & \textbf{解析解 (rad/s)} & \textbf{誤差 (\%)} \\
\midrule
20 & 20184.54 & 19992.40 & 0.9611 \\
30 & 19989.48 & 19992.40 & 0.0146 \\
40 & 19921.78 & 19992.40 & 0.3532 \\
50 & 19890.54 & 19992.40 & 0.5095 \\
100 & 19848.99 & 19992.40 & 0.7173 \\
\bottomrule
\end{tabular}
\end{center}

\subparagraph{Mode 4 收斂性分析：}

\begin{center}
\begin{tabular}{cccc}
\toprule
\textbf{元素數} & \textbf{FEM (rad/s)} & \textbf{解析解 (rad/s)} & \textbf{誤差 (\%)} \\
\midrule
20 & 39990.04 & 39176.41 & 2.0768 \\
30 & 39217.83 & 39176.41 & 0.1057 \\
40 & 38952.32 & 39176.41 & 0.5720 \\
50 & 38830.25 & 39176.41 & 0.8836 \\
100 & 38668.30 & 39176.41 & 1.2970 \\
\bottomrule
\end{tabular}
\end{center}

\textit{注：震動頻率誤差不隨元素增加而持續下降，這是因為 Timoshenko 理論考慮剪力變形與旋轉慣量，導致頻率略低於 Euler-Bernoulli 解析解。此差異為理論本身固有特性，非網格細化所能消除。高模態（Mode 3、4）在較少元素時誤差較大，需要較細的網格才能收斂。}

\section{結論}

本報告成功實現了 Timoshenko 樑的有限元素分析，並通過與解析解的對比驗證了程式碼的正確性。主要結論如下：

\begin{enumerate}
  \item \textbf{靜態彎曲分析：} FEM 結果與解析解高度吻合，誤差不超過 0.4\%，展現良好的收斂性。
  
  \item \textbf{屈曲分析：} 臨界載重計算結果穩定收斂至解析解，驗證了几何剛度矩陣實現的正確性。
  
  \item \textbf{自由震動分析：} 低階模態的頻率預測準確，高階模態與 Euler-Bernoulli 理論存在固有差異，這正是 Timoshenko 理論考慮剪力變形與旋轉慣量的體現。
  
  \item \textbf{數值技巧：} 採用單點積分有效消除剪力鎖定問題，保證了薄樑分析的準確性。
\end{enumerate}

本專案開發的 Julia 程式碼可作為 Timoshenko 樑有限元素分析的參考實現，適用於工程實踐與學術研究。

\end{document}
